\chapter{Statik}

In der Statik - als Teilgebiet der Mechanik - wird von ruhenden und unverformbaren 
(starren) Körpern ausgegangen. Ziel der Statik ist die Annahme der Lasten an einem
Bauteil und die Bestimmung von Lagerreaktionen.

Diese Annahmen werden benötigt, um die Festigkeit in einem Bauteil zu bestimmen.
Festigkeit ist nicht Teil dieses Kapitels.
Abbildung \ref{fig:übersichtMechanikStatik} zeigt
eine Einteilung der Teilgebiete in der Mechanik.

\begin{figure}[H]    
    \centering          %     left botm right top
    \includegraphics[clip, trim=0cm 0cm 0cm 0cm,width=.8\textwidth]{übersichtMechanikStatik.png} %.6 stands for 0.6 x textwidth
    \caption{Einteilung Mechanik und Statik}
    \label{fig:übersichtMechanikStatik}
\end{figure}

Die Berechnungen und Beispiele in diesem Kapitel werden hauptsächlich in der Ebene (2D)
gezeigt, wobei die meisten Konzepte in den Raum übertragen werden können (3D).

Viele Inhalte, Bilder und Aufgaben dieses Kapitels stammen aus dem Mechanik \&
Festigkeit Skript von Prof. R. Bärtsch, unterrichtet an der HSLU T\&A in Luzern 
\cite{bartsch_mechanik_2022}.

\section{Kraft}
Eine Kraft selbst ist nicht sichtbar, deren Wirkung kann jedoch 
sichtbar sein. Sie kann beispielsweise eine Masse beschleunigen
oder Deformationen verursachen. Die Kraft ist die zentrale Grösse
in der Statik.

3 Angaben sind zur Beschreibung einer Kraft notwendig:

\begin{minipage}{.5\textwidth}
    % Enumerate environment on the left
    \begin{itemize}
        \item Der Betrag
        \item Die Wirkungslinie
        \item Der Richtungssinn
    \end{itemize}
\end{minipage}%
\begin{minipage}{.5\textwidth}
    % Image on the right
    \centering
    \captionsetup{type=figure}
    \includegraphics[clip, trim=0cm 0cm 0cm 0cm, width=\textwidth]{kraftvektor.png}
    \captionof{figure}{Kraftvektor}
    \label{fig:kraftvektor}
\end{minipage}

Die Kraft ist also eine vektorielle Grösse und wird mit einem Pfeil dargestellt. Die ebene
Kraft $F$ kann in ihre Komponenten $F_x$ und $F_y$ zerlegt werden
(Abbildung \ref{fig:kraftImKartesischen}).

\begin{minipage}{.5\textwidth}
    \begin{equbox}
        \begin{equation}
            \begin{split}
                F_x = F \cdot \cos \alpha \\
                F_y = F \cdot \sin \alpha \\
                \tan \alpha = \frac{F_y}{F_x} \\
                F = \sqrt{F_x^2 + F_y^2}
            \end{split}
            \label{eq:equationKraft}
        \end{equation}
    \end{equbox}
\end{minipage}%
\begin{minipage}{.5\textwidth}
    \centering
    \captionsetup{type=figure}
    \includegraphics[clip, trim=0cm 0cm 0cm 0cm, height=4cm]{kraftImKartesischen.png}
    \captionof{figure}{Resultierende Kraft}
    \label{fig:kraftImKartesischen}
\end{minipage}

Die Wirkung der Kraft $F$ mit Wirkungslinie in Richtung von $\alpha$ ist die gleiche wie diejenige
der Kräfte $F_x$ und $F_y$. Man sagt, die beiden Kräftegruppen sind äquivalent.

\subsection{Die resultierende Kraft}

Zwei Kräfte, die an einem Punkt angreifen, können durch eine resultierende Kraft ersetzt werden. 
Die Erweiterung auf mehrere Kräfte mit beliebigen Angriffspunkten führt auf die 
Definition der Resultierenden.

\begin{minipage}{.5\textwidth}
    \begin{infobox}
        Eine Kraft, die die gleiche Wirkung hat wie eine Reihe von Kräften, 
        die sie ersetzt, nennt man die Resultierende.
    \end{infobox}
\end{minipage}%
\begin{minipage}{.5\textwidth}
    \centering
    \captionsetup{type=figure}
    \includegraphics[clip, trim=0cm 0cm 0cm 0cm, width=.7\textwidth]{resultierende Kraft.png}
    \captionof{figure}{Ersatz von Kräften}
    \label{fig:resultierende}
\end{minipage}

\section{Moment}

Eine Kraft F im Abstand d von einer Achse bewirkt eine Drehung des Körpers, wobei die
Wirkung proportional zur Kraft F und zum Abstand d ist.

\begin{equbox}
    \begin{equation}
        M = F \cdot d, F \perp d
    \end{equation}
\end{equbox}

\begin{figure}[H]    
    \centering          %     left botm right top
    \includegraphics[clip, trim=0cm 0cm 0cm 0cm,width=\textwidth]{Das Moment einer Kraft.png} %.6 stands for 0.6 x textwidth
    \caption{Das Moment einer Kraft}
    \label{fig:dasMomentEinerKraft}
\end{figure}

Bei der Berechnung des Momentes einer Kraft bezüglich eines Punktes A in der Ebene kann
entweder gemäss Definition mit dem senkrechten Abstand $d$ der Wirkungslinie zum Pol
gerechnet werden. Alternativ - und vielfach einfacher - kann die Kraft zunächst in ihre
Komponenten zerlegt und dann die Einzelmomente dieser beiden Kraftkomponenten
berechnet und anschliessend summiert werden.

\begin{infobox}
    Das resultierende Moment mehrerer Kräfte ist gleich der Summe der Einzelmomente.
\end{infobox}

Da sich eine Kraft entlang ihrer Wirkungslinie verschieben lässt, ohne dass ihre Wirkung sich
ändert, gibt es eine spezielle Lage der Kraft $F$, für die die Komponenten $F_x$ durch den Pol geht
und damit kein Moment mehr verursacht. Alle drei Varianten in Abbildung 
\ref{fig:zurBerechnungDesMomentsEinerKraft} ergeben das gleiche Ergebnis.

\begin{figure}[H]    
    \centering          %     left botm right top
    \includegraphics[clip, trim=0cm 0cm 0cm 0cm,width=\textwidth]{zurBerechnungDesMomentsEinerKraft.png} %.6 stands for 0.6 x textwidth
    \caption{Zur Berechnung des Momentes einer Kraft}
    \label{fig:zurBerechnungDesMomentsEinerKraft}
\end{figure}

\subsection{Das Kräftepaar}

Das Kräftepaar hat die Wirkung eines Momentes, dessen Grösse F·d unabhängig von der
Lage des Pols ist. Daraus ergibt sich folgende Erkenntnis:

\begin{infobox}
    Ein Kräftepaar (oder Moment) kann in der Statik beliebig verschoben werden, ohne dass sich
    seine Wirkung ändert.
\end{infobox}

\begin{figure}[H]    
    \centering          %     left botm right top
    \includegraphics[clip, trim=0cm 0cm 0cm 0cm,width=\textwidth]{verschiebenVonKräftepaaren.png} %.6 stands for 0.6 x textwidth
    \caption{Verschieben von Kräftepaaren}
    \label{fig:verschiebenVonKräftepaaren}
\end{figure}

\subsection{Schwerpunkt}

Der Schwerpunkt befindet sich dort, wo man seinen Zeigefinger an einem Körper positionieren müsste, 
um es im Gleichgewicht zu halten bzw. um diesen zu balancieren \cite{noauthor_schwerpunkt_nodate}.
In der Statik ist man in erster Linie an der Lage des Schwerpunktes interessiert 
($S$ in Abbildung \ref{fig:lageSchwerpunktAllgemein}).

\begin{figure}[H]    
    \centering          %     left botm right top
    \includegraphics[clip, trim=0cm 0cm 0cm 0cm,width=.5\textwidth]{schwerpunkt_allgemein.png} %.6 stands for 0.6 x textwidth
    \caption{Lage des Schwerpunktes an einer allgemeinen Fläche \cite{noauthor_schwerpunkt_nodate}}
    \label{fig:lageSchwerpunktAllgemein}
\end{figure}

Nachfolgend die allgemeine
Formel für die Lage in $x$ und $y$. Selbstverständlich kann dieses Konzept auf einen
Körper (3D) erweitert werden.

\begin{equbox}
    \begin{align}
        x_s &= \frac{\sum A_i \cdot x_{Si}}{\sum A_i} \\
        y_s &= \frac{\sum A_i \cdot y_{Si}}{\sum A_i}
    \end{align}
\end{equbox}

Um den Schwerpunkt zu berechnen, kann in einem ersten Schritt eine Grundfläche definiert werden.
Zu dieser werden anschliessend Unterflächen (gemäss obigen Formeln) addiert oder subtrahiert, 
wie in Abbildung \ref{fig:schwerpunktberechung} gezeigt.

$x_{Si}$ und $y_{Si}$ sind dabei die Positionen der Schwerpunkte der Unterflächen.

\begin{figure}[H]    
    \centering          %     left botm right top
    \includegraphics[clip, trim=0cm 0cm 0cm 0cm,width=.8\textwidth]{schwerpunkt_übersicht.png} %.6 stands for 0.6 x textwidth
    \caption{Konzept zur Schwerpunktberechnung \cite{noauthor_schwerpunkt_nodate}}
    \label{fig:schwerpunktberechung}
\end{figure}

Für komplexere Berechnungen kann eine Tabelle zur Hilfe verwendet werden.

\begin{table}[h]
\centering
\caption{Vorlage Tabelle für Berechnungen}
\label{tab:leere_tabelle}
\begin{tabular}{|c|c|c|c|c|c|}
\hline
$i$ & $A_i$ & $x_{Si}$ & $y_{Si}$ & $A_i \cdot x_{Si}$ & $A_i \cdot y_{Si}$ \\ \hline
1 & & & & & \\ \hline
2 & & & & & \\ \hline
3 & & & & & \\ \hline
$\sum$ & & - & - & & \\ \hline
\end{tabular}
\end{table}

\section{Gleichgewicht eines Systems}

Die Berechnung der Resultierenden dient der Ermittlung von Kräften an einem Bauteil, das
sich im Gleichgewicht befindet. Gleichgewicht bedeutet, dass sich alle angreifenden Kräfte in
ihrer Wirkung aufheben, also die Resultierende verschwindet.

Die Bestimmung der Resultierenden wird schrittweise für die folgenden Arten von Kräfte- systemen gelöst.

\begin{itemize}
    \item Zentrales Kräftesystem (gemeinsamer Angriffspunkt)
    \item System mit parallelen Kräften (parallele Wirkungslinien)
    \item Allgemeines Kräftesystem (beliebige Angriffspunkte und Wirkungslinien)
\end{itemize}

\begin{figure}[H]    
    \centering          %     left botm right top
    \includegraphics[clip, trim=0cm 0cm 0cm 0cm,width=\textwidth]{die dre Arten von kräftesystemen.png} %.6 stands for 0.6 x textwidth
    \caption{Die drei Arten von Kräftesystemen}
    \label{fig:dieDreiArten}
\end{figure}

\subsection{Zentrales Kräftesystem}

Das zentrale Kräftesystem repräsentiert ein Punkt, an dem beliebig viele Kräfte
angreifen. Aus der Summe der Kräfte kann eine Resultierende Kraft berechnet werden,
gezeigt in Abbildung \ref{fig:kräfteplan} und Formeln \ref{eq:komAddition} bis
\ref{eq:RichtungDerResultierenden}.

\begin{figure}[H]    
    \centering          %     left botm right top
    \includegraphics[clip, trim=0cm 0cm 0cm 0cm,width=\textwidth]{kräfteplan.png} %.6 stands for 0.6 x textwidth
    \caption{Resultierende Kraft im Kräfteplan}
    \label{fig:kräfteplan}
\end{figure}

Anhand des Kräfteplans kann eine graphische Lösung ermittelt werden, er ist aber
nicht notwendig für die rechnerische Lösung. Dazu genügt der Lageplan.

\begin{equbox}
    \begin{minipage}[t]{.25\textwidth}
        Komponentenweise \newline Addition:
        \begin{equation}
            \begin{split}
                &F_{res, x} = \sum F_x \\
                &F_{res, y} = \sum F_y
            \end{split}
            \label{eq:komAddition}
        \end{equation}
    \end{minipage}%
    \hspace{.05\textwidth}
    \begin{minipage}[t]{.34\textwidth}
            Der Betrag der Resultierenden:
            \begin{equation}
                F_{res} = \sqrt{F_{res, x}^2 + F_{res, y}^2}
                \label{eq:betragDerResultierende}
            \end{equation}
    \end{minipage}
    \hspace{.05\textwidth}
    \begin{minipage}[t]{.25\textwidth}
            Die Richtung der \newline Resultierenden:
            \begin{equation}
                \tan{\alpha} = \frac{F_{res, y}}{F_{res, x}}
                \label{eq:RichtungDerResultierenden}
            \end{equation}
    \end{minipage}
\end{equbox}

\begin{bspbox}
    \begin{minipage}{.66\textwidth}
        3 Kräfte greifen an einem Haken an, gesucht ist die resultierende Kraft nach
        Betrag und Richtung. \newline
        $\displaystyle F_{res, x} = 6kN \cdot \cos{30^{\circ}} + 8kN \cdot \cos{60^{\circ}}
        - 5kN \cdot \sin{20^{\circ}} = 7.49kN$ \newline
        $\displaystyle F_{res, y} = 6kN \cdot \sin{30^{\circ}} + 8kN \cdot \sin{60^{\circ}}
        + 5kN \cdot \cos{20^{\circ}} = 14.63kN$ \newline
        $\displaystyle F_{res} = \sqrt{(7.49kN)^2 + (14.63kN)^2} = 16.4kN$ \newline
        $\displaystyle \tan{\alpha} = \frac{14.63kN}{7.49kN} = 1.953 \to \alpha = 62.9^{\circ}$ 
    \end{minipage}
    \hspace{.04\textwidth}
    \begin{minipage}{.28\textwidth}
        \begin{figure}[H]    
            \centering          %     left botm right top
            \includegraphics[clip, trim=0cm 0cm 0cm 0cm,width=\textwidth]{zentralesKräftesystemHakenBeispiel.png} %.6 stands for 0.6 x textwidth
            \label{fig:bspHaken}
        \end{figure}
    \end{minipage}
\end{bspbox}

Zu beachten ist die Mehrdeutigkeit der $arctan$ Funktion. Zur Kontrolle der Richtung der
Resultierenden kann eine grobmassstäbliche Skizze des Kräfteplans dienen.

\subsection{Parallele Kräfte}

Die Resultierende von parallelen Kräften muss die gleiche Wirkung wie die Einzelkräfte
haben. Daraus ergeben sich folgende Bedingungen:

\begin{infobox}
    \begin{itemize}
        \item Der Betrag der Resultierenden Kraft ist gleich der Summe der Einzelkräfte.
        \item Das Moment der Resultierenden ist gleich der Summe der Momente aus den Einzelkräften.
        \begin{itemize}
            \item Diese Bedingung muss für jeden beliebigen Pol erfüllt sein.
        \end{itemize}
    \end{itemize}
\end{infobox}

\begin{minipage}{.5\textwidth}
    \begin{equbox}
        \begin{equation}
            \begin{split}
                F_{res} = \sum F_i \\
                x_s = \frac{\sum F_i \cdot x_i}{\sum F_i}
            \end{split}
            \label{eq:resParallel}
        \end{equation}
    \end{equbox}
\end{minipage}%
\begin{minipage}{.5\textwidth}
    \centering
    \captionsetup{type=figure}
    \includegraphics[clip, trim=0cm 0cm 0cm 0cm, height=5cm]{resultierendeParalleleKräfte.png}
    \captionof{figure}{Resultierende parall. Kräfte}
    \label{fig:resParallel}
\end{minipage}

\begin{bspbox}
    \begin{minipage}{.58\textwidth}
        Der Betrag der Resultierenden ist: \newline
        $\displaystyle F_{res} = F_1 + F_2 + F_3$ \newline
        Das Moment bezüglich des Pols in A ist: \newline
        $\displaystyle F_{res} \cdot x_s = F_1 \cdot x_1 + F_2 \cdot x_2 + F_3 \cdot x_3$ \newline
        $\displaystyle x_s = \frac{F_1 \cdot x_1 + F_2 \cdot x_2 + F_3 \cdot x_3}{F_{res}}$
    \end{minipage}
    \begin{minipage}{.4\textwidth}
        \begin{figure}[H]   
            \centering          %     left botm right top
            \includegraphics[clip, trim=0cm 0cm 0cm 0cm,width=\textwidth]{beispielParalleleKräfte.png} %.6 stands for 0.6 x textwidth
            \label{fig:bspParallel}
        \end{figure}
    \end{minipage}
\end{bspbox}

\subsection{Allgemeines Kräftesystem}

Bei allgemeinen Kräftesystemen sind die Kräfte weder parallel noch haben sie einen
gemeinsamen Schnittpunkt. Die Resultierende muss auch hier die gleiche Wirkung wie die
Einzelkräfte haben.

\begin{minipage}{.49\textwidth}
    \begin{infobox}
        Ein allgemeines ebenes Kräftesystem kann im beliebigen Punkt der Ebene durch eine
        resultierende Kraft und ein Moment ersetzt werden. Das Moment ist dabei nicht an diesen
        Punkt gebunden, sondern kann in der Ebene frei verschoben werden.
    \end{infobox}
\end{minipage}%
\hspace{.01\textwidth}
\begin{minipage}{.49\textwidth}
    \begin{equbox}
        Formal sind also die 3 Bedingungen zu erfüllen:
        \begin{equation}
            \begin{split}
                &F_{res, x} = \sum F_x \\
                &F_{res, y} = \sum F_y \\
                &M_{res} = \sum M
            \end{split}
        \end{equation}
    \end{equbox}
\end{minipage}


\begin{bspbox}
    \begin{minipage}{.58\textwidth}
        Auf das Bauteil wirken die vier Kräfte:
        $F_1 = 100N, F_2 = 50N, F_3 = 200N$ und $F_4 = 150N$.
        Zu bestimmen ist die resultierende Kraft nach Lage, Grösse
        und Richtung.
        Die resultierende Kraft ist:
        \begin{equation*}
            \begin{split}
                &F_{res, x} = F_1 + F_3 \cdot \cos{45^{\circ}} = 241.4N \\
                &\displaystyle F_{res, y} = -F_2 - F_3 \cdot \sin{45^{\circ}} - F_4 = -341.4N \\
                &\displaystyle F_{res} = \sqrt{(241.1N)^2 + (-341.4N)^2} = 418.1N
            \end{split}
        \end{equation*}
        Das resultierende Moment bez. Punkt A ist:
        \begin{equation*}
            \begin{split}
                &M_A = -F_2 \cdot 0.1m - F_3 \cdot \sin{45^{\circ}} \cdot 0.2m \\
                &- F_4 \cdot 0.2m = -63.3Nm
            \end{split}
        \end{equation*}
        Die Richtung der resultierenden Kraft is:
        \begin{equation*}
            \begin{split}
                &\tan{\alpha} = \frac{-341.4N}{241.4N} = -1.414 \\
                &\rightarrow \alpha = -54.7^{\circ}
            \end{split}
        \end{equation*}
        Um die Lage der resultierenden Kraft zu bestimmen,
        wird die Resultierende so lange nach rechts
        verschoben, bis das Moment der resultierenden Kraft
        um A gleich gross wie $M_A$ wird, d.h. $F_{res, y} \cdot x = M_A$.
        Die Komponente $F_{res, x}$ ergibt kein Moment bez. A.
        \begin{equation*}
            x = \frac{|M_A|}{|F_{res, y}|} = \frac{|-63.3Nm|}{|-341.4N|} = 0.1854m
        \end{equation*}
        Analog könnte man die Resultierende auch nach oben
        verschieben und den y-Abstand berechnen.
    \end{minipage}
    \begin{minipage}{.4\textwidth}
        \begin{figure}[H]
            \centering
            \begin{subfigure}{\textwidth}
                \includegraphics[clip, trim=0cm 0cm 0cm 0cm, width=\textwidth]{beispielAllgKräftesys01.png}
                \label{fig:beispielAllgKräftesys01}
            \end{subfigure}
            \begin{subfigure}{\textwidth}
                \includegraphics[clip, trim=0cm 0cm 0cm 0cm, width=\textwidth]{beispielAllgKräftesys02.png}
                \label{fig:beispielAllgKräftesys02}
            \end{subfigure}
            \begin{subfigure}{\textwidth}
                \includegraphics[clip, trim=0cm 0cm 0cm 0cm, width=\textwidth]{beispielAllgKräftesys03.png}
                \label{fig:beispielAllgKräftesys03}
            \end{subfigure}
        \end{figure}
    \end{minipage}
\end{bspbox}

\section{Lagerung und Freimachen}

Fast alle Bauteile sind in ihrer Umgebung gelagert, z.B. an einem Fundament festgeschraubt
oder in einem Gelenk drehbar gelagert. Ausnahmen sind schwimmende oder fliegende
Bauteile. Die Lagerung hat die Aufgabe, das Bauteil zu positionieren, Kräfte zu übertragen
und Bewegungsfreiheitsgrade einzuschränken.
Die Aufgabe der Statik ist es, diese Lagerkräfte zu berechnen. Es muss deshalb zunächst
Klarheit bestehen, wo und wie diese wirken. Die Lager und Verbindungselemente übertragen
je nach konstruktiver Ausführung verschiedene Belastungen. Um zu erkennen, welche Kräfte
die Lagerstellen auf ein Bauteil ausüben, wird dieses aus seiner Verbindung gelöst, d.h.
freigemacht.

Das Freimachen geschieht durch einen (gedachten) Schnitt, welcher das Bauteil von den
Lagern trennt. An den geschnittenen Stellen werden Kräfte bzw. Momente eingeführt, die
sog. Lagerreaktionen. Sowohl die Lagerreaktionen als auch die wirkenden Lasten sind
äussere Kräfte, da sie von aussen auf das Bauteil einwirken. Das freigemachte Bauteil ist in
Ruhe, d.h. alle an ihm angreifenden Kräfte und Momente sind im Gleichgewicht. Dies erlaubt
es, die Lagerreaktionen zu berechnen \cite{bartsch_statik_2022}.

\subsection{Das Freimachen}

Der dargestellte Generator (Abbildung \ref{fig:freimachengenerator}) ist in A in einem Gelenk gelagert und wird in B von einem Stab
gehalten. Die Belastung besteht aus den Kräften F und G und dem Moment M. Der gedachte
Schnitt trennt die beiden Lagerelemente. Würde man den Stab durchtrennen, so müsste man
den Generator mit einer Kraft B halten, welche genau so gross ist, wie die Stabkraft vor dem
Trennen. Anstelle des Gelenkes A wird eine Kraft A in unbekannter Richtung und Grösse
eingeführt \cite{bartsch_statik_2022}.

\begin{figure}[H]    
    \centering          %     left botm right top
    \includegraphics[clip, trim=0cm 0cm 0cm 0cm,width=.7\textwidth]{freimachen_generator.png} %.6 stands for 0.6 x textwidth
    \caption{Freimachen eines belasteten Bauteils \cite{bartsch_statik_2022}}
    \label{fig:freimachengenerator}
\end{figure}

Das Freimachen beinhaltet:

\begin{itemize}
    \item Skizzieren des Bauteils (vereinfacht), herausgelöst aus allen Verbindungen
    \item Übertragen aller äusseren Kräfte / Momente, die auf das Bauteil wirken
    \item Einführen und Bezeichnen der Kräfte / Momente an den geschnittenen Stellen
\end{itemize}

\subsection{Lagerungsarten}

Bild \ref{fig:lagerungsartenUndLagerreaktionenEbene} zeigt häufig vorkommende Lagerungsarten mit den zugehörigen Freiheitsgraden

\begin{figure}[H]    
    \centering          %     left botm right top
    \includegraphics[clip, trim=0cm 0cm 0cm 0cm,width=\textwidth]{lagerungsarten.png} %.6 stands for 0.6 x textwidth
    \caption{Lagerungsarten und Lagerreaktionen in der Ebene \cite{bartsch_statik_2022}}
    \label{fig:lagerungsartenUndLagerreaktionenEbene}
\end{figure}

\begin{bspbox}
    Am freigemachten Bauteil ist der Wirkungssinn der unbekannten Lagerkräfte grundsätzlich
    belanglos und damit frei wählbar. Den richtigen Wirkungssinn ergibt das Vorzeichen der
    Berechnung. Ein positives Vorzeichen bedeutet, Kraft bzw. Moment wirken wie
    angenommen; ein negatives Vorzeichen, Kraft bzw. Moment wirken umgekehrt wie
    angenommen. Wenn es die Anschauung erlaubt, ist es sinnvoll, die Lagerreaktionen in der
    „richtigen“ Richtung einzuzeichnen.

    \begin{figure}[H]    
        \centering          %     left botm right top
        \includegraphics[clip, trim=0cm 0cm 0cm 0cm,width=\textwidth]{beispiel_freimachen.png} %.6 stands for 0.6 x textwidth
        \label{fig:bspfreimachen}
    \end{figure}
\end{bspbox}


